\exnumber{02}
\extitle{UI IV - Ex 3 - Écrans}
\exfiles{ex02/fichier Figma, nomme Name\_FirstName\_DeliverableName\_Date\_VersionNumber}
\exauthorize{Ceux nommés dans l'exercice}
\makeheaderfiles

\textbf{Instructions :}

À partir des wireframes et de votre library de design system, amenez votre prototype moyenne fidélité à un niveau haute fidélité !

Sélectionnez 5 écrans qui montreront le user flow d'un client qui :

\begin{itemize}
\item arrive sur la page d'accueil
\item parcourt le contenu
\item consulte une page produit
\item ajoute le produit au panier
\item finalise l'achat
\end{itemize}

Utilisez vos composants, call to actions et visuels pour compléter les wireframes.

\textbf{Livrables :}

\begin{itemize}
\item 5 écrans haute fidélité depuis Figma
\item Déposez le fichier Figma sur Github, nommé ainsi : « Name\_FirstName\_DeliverableName\_Date\_VersionNumber ». Exemple : « Doe\_John\_MidFiPrototype\_05242024\_V1 ».
\item Ajoutez un titre dans le fichier.
\item Soignez l'aspect global du fichier. Il doit être clair, simple et lisible. Vous pouvez regarder canva.com, sans surcharger le dessin de détails.
\end{itemize}

\textbf{Ressources :}

\begin{itemize}
\item Unsplash \url{https://unsplash.com/} pour les visuels
\item De la basse fidélité à la haute fidélité \url{https://mockitt.wondershare.com/prototyping/low-fidelity-vs-high-fidelity.html}
\end{itemize}

